\section{The process, its roots, and its closed forms}
\label{dist:sec:recap}

The birth--death--catastrophe (BDC) process was defined and solved in
\ChCore. Four things are needed from it here: the notation, the process and
its release counter, the roots, and the identity relating them.
\Cref{dist:app:formulae} collects the closed forms.

\subsection{Canonical notation}
\label{dist:sec:notation-tab}

\Cref{dist:tab:notation} fixes the notation, taken unchanged from \ChCore.
Every symbol appearing in more than one section is listed; those local to a
single derivation are introduced where they occur.

\begin{table}[tbp]
\centering
\renewcommand{\arraystretch}{1.05}
\footnotesize
\begin{tabular}{@{}llp{0.40\textwidth}@{}}
\toprule
Symbol & Meaning & Notes \\
\midrule
$\lambda,\mu,\delta$ & birth, death, catastrophe rates & per capita \\
$\Xt$ & intracellular load & a CTMC on $\{0,1,2,\ldots\}\cup\{H\}$ \\
$H$ & absorbing rupture state & \\
$\Wt$ & released count from one cell & $0$ until the burst \\
$\tau$ & burst time & possibly $\infty$ \\
$\Kb$ & burst size & never bare $K$, since $K(t)=\E{\Xt^2}$ \\
$G(z,t)$ & probability generating function (PGF) of the load & defective under catastrophe \\
$G_k$ & PGF from $k$ founders & $=G^k$ \\
$p_n(t)$ & $\pr{\Xt=n}$ & \\
$P(t)$, $\sigma$ & geometric ratio of the load; the clock $\sigma=P(t)$ & slides from $0$ to $1/a$ \\
$\xi$, $s$ & characteristic coordinate and parameter & $t=-\xi$; \cref{dist:sec:pgf}. In \cref{dist:sec:chained}, $s=\delta/(\lambda+\delta)$ \\
$I(t)$, $D(t)$ & no burst by $t$; cleared internally, no burst & $D=p_0(t)$ \\
$\Ifix(t)$ & $\pr{\Xt\notin\{0,H\}}=I-D$ & productive infection \\
$J(t),K(t)$ & $\E{\Xt}$, $\E{\Xt^2}$ & \\
$V(t)$ & $\E{\Wt}$ & cumulative mean release \\
$\varphi(t)$ & density of $\tau$ & defective, mass $1-b$ \\
$a,b$ & characteristic roots & $b<1<a$ \\
$A,B,\theta$ & $a-1$, $1-b$, $\lambda(a-b)$ & $AB=\delta/\lambda$ \\
$\kappa$, $w$ & $1+\delta/(2\lambda)$; \ $\ee^{\theta t}$ & in $K(t)$; the recurring exponential \\
$q_n,\nu$ & quasi-stationary law & $\nu$ when read as a vector \\
$\vartheta$, $\mathcal Q$ & decay rate; killed subgenerator & $\vartheta=\theta$; \cref{dist:prop:decay} \\
$T_{\rm prod}$, $H_n$ & productive lifetime; harmonic number $\sum_{j\le n}1/j$ & \\
$g_k(t)$ & release flux from $k$ founders & $=\delta K_k(t)$ \\
$r_k,t_k,T_k$ & size, time and interval of the $k$th rupture & \cref{dist:sec:chained} \\
$\rho$ & $\lambda/(\lambda+\delta)=1-s$ & \cref{dist:sec:chained} \\
$d_{\Icell},p$ & budding death and release rates & \cref{dist:sec:budding} \\
$\Fhyp$ & Gauss hypergeometric function & \\
\bottomrule
\end{tabular}
\caption{Canonical notation for this chapter.}
\label{dist:tab:notation}
\end{table}

\subsection{The process and its release counter}
\label{dist:sec:process}

The object throughout is one cell holding one intracellular population.

\begin{definition}[Birth--death--catastrophe process]
\label{dist:def:bdc}
The load $\Xt$ is a continuous-time Markov chain on
$\{0,1,2,\ldots\}\cup\{H\}$ which, from any state $n\ge1$, jumps
\begin{align*}
n &\to n+1 &&\text{at rate } \lambda n \quad\text{(birth)},\\
n &\to n-1 &&\text{at rate } \mu n \quad\text{(death)},\\
n &\to H &&\text{at rate } \delta n \quad\text{(catastrophe)}.
\end{align*}
Both $0$ and $H$ are absorbing: the population either clears itself
internally, or the cell ruptures. Unless stated otherwise $X_0=1$.
\end{definition}

The classical treatments
\cite{brockwell1982birth,karlin1982linear,di2008note} allow immigration and
distributed catastrophe sizes; this specialisation is the one with the
closed-form family we need. Rupture is what matters, so the release count runs
alongside the load: with $\tau$ the catastrophe time,
\begin{equation}
\Wt = \begin{cases}
0, & t<\tau,\\
X_{\tau^-}, & t\ge\tau,
\end{cases}
\label{dist:eq:wdef}
\end{equation}
so that at rupture the whole pre-burst load is dumped and then held.
\Cref{dist:fig:joint} shows one path of the pair; every single-cell release
quantity below is a functional of $(\Xt,\Wt)$.

\begin{figure}[tbp]
\centering
\includegraphics[width=0.66\textwidth]{fig01_joint_process/fig01}
\caption{A representative path of the joint process $(\Xt,\Wt)$. The
intracellular load follows a birth--death path until rupture at time $\tau$,
when $\Xt$ enters the absorbing state $H$. At the same instant the released
count jumps from zero to the pre-rupture load $X_{\tau^-}$ and remains there.
The two axes share a time origin: nothing leaves the cell before $\tau$, and
everything leaves at $\tau$.}
\label{dist:fig:joint}
\end{figure}

\subsection{Roots and the workhorse identity}
\label{dist:sec:roots}

Write
\begin{equation}
\eta=\frac{\lambda+\mu+\delta}{2\lambda},\qquad
a=\eta+\sqrt{\eta^2-\tfrac{\mu}{\lambda}},\qquad
b=\eta-\sqrt{\eta^2-\tfrac{\mu}{\lambda}},
\label{dist:eq:roots}
\end{equation}
the two roots of $\lambda f^2-(\lambda+\mu+\delta)f+\mu$, and set
\begin{equation}
A:=a-1>0,\qquad B:=1-b>0,\qquad \theta:=\lambda(a-b)
=\sqrt{(\lambda+\mu+\delta)^2-4\lambda\mu}.
\label{dist:eq:abtheta}
\end{equation}
Vieta's relations give $ab=\mu/\lambda$ and
$a+b=(\lambda+\mu+\delta)/\lambda$, and the identity doing most of the
algebraic work below is
\begin{equation}
AB=(a-1)(1-b)=(a+b)-ab-1=\frac{\delta}{\lambda},
\qquad A+B=a-b.
\label{dist:eq:ab}
\end{equation}
That $b<1<a$ follows immediately from
$\lambda(1-a)(1-b)=\lambda(1-(a+b)+ab)=-\delta<0$, which places $1$ strictly
between the roots.

\fwd{twotype}{\ChTwoTypecap} solves a Riccati equation of exactly this
shape under unrelated names: roots that straddle one, a root gap playing the
part of $a-b$, and the same factorisation. \Cref{dist:sec:next} records the
correspondence rather than acting on it.

\subsection{The closed forms}
\label{dist:sec:moments-recap}

There are three absorbing-behaviour probabilities, not two.
\begin{definition}[Fixation functions]
\label{dist:def:fixation}
\begin{equation}
I(t)=\pr{\text{no burst by }t},\qquad
D(t)=\pr{\text{internal extinction by }t,\ \text{no burst}},
\label{dist:eq:idefs}
\end{equation}
\begin{equation}
\Ifix(t)=\pr{\Xt\notin\{0,H\}}=I(t)-D(t),
\label{dist:eq:ifixdef}
\end{equation}
the last being the probability that the cell is still productively infected.
\end{definition}

The first two satisfy the same Riccati equation,
\begin{equation}
\ddt{I}=\lambda(I-a)(I-b)=\mu-(\lambda+\mu+\delta)I+\lambda I^2,
\label{dist:eq:riccati}
\end{equation}
differently initialised ($I(0)=1$, $D(0)=0$); its branching term $I^2$ comes
from the two independent daughter lineages that a birth of the founder
creates.
Separating variables, with $w:=\ee^{\theta t}\ge1$, gives the triplet
\begin{equation}
\boxed{I(t)=\frac{aB+bAw}{B+Aw},\qquad
D(t)=\frac{ab\,(w-1)}{aw-b},\qquad
\Ifix(t)=\frac{(a-b)^2\,w}{(B+Aw)(aw-b)},}
\label{dist:eq:idihat}
\end{equation}
with $I(\infty)=D(\infty)=b$, $\Ifix(\infty)=0$, $\Ifix(0)=1$ and
$\Ifix'(0)=-(\mu+\delta)$: at founding the productive state is lost only by
the death or the catastrophe of the single founder. \Cref{dist:cor:sum}
confirms the third formula independently.

The remaining closed forms are derived in \ChCore. The hazard of bursting is
$\delta$ times the current load, so $I'=-\delta J$, and
\begin{align}
J(t)&=\E{\Xt}=-\frac1\delta\ddt{I}=\frac{(a-b)^2w}{(B+Aw)^2}
=-\frac{\lambda}{\delta}(I-a)(I-b),
\label{dist:eq:j}\\
K(t)&=\E{\Xt^2}=\Bigl[1+\tfrac{2\lambda}{\delta}(1-I)\Bigr]J
=\frac{2\lambda}{\delta}(\kappa-I)J,\qquad \kappa=1+\frac{\delta}{2\lambda},
\label{dist:eq:k}\\
V(t)&=\E{\Wt},\qquad V'(t)=\delta K(t),\qquad
V(t)=(1-I)\Bigl[1+\tfrac{\lambda}{\delta}(1-I)\Bigr],
\label{dist:eq:v}\\
V_\infty&=\frac{a(1-b)}{a-1}=\frac{aB}{A},
\qquad
\E{\Kb\mid\text{burst}}=\frac{V_\infty}{1-b}=\frac{a}{a-1}.
\label{dist:eq:vinf}
\end{align}
One small point carries large consequences.
$V_\infty=\E{\Kb\ind\{\text{burst}\}}$ \emph{includes} the realisations that
cleared internally and released nothing, a mass $b$ at $0$, whereas the
conditional mean removes them; \cref{dist:sec:cond-means} takes the
distinction apart.

The assay predictions of \cref{dist:sec:discussion} need the second moment of
the release as well. \ChCorecap{} obtains it from the moment hierarchy,
eliminating $\E{\Xt^3}$ through the equation for $K$:
\begin{equation}
\boxed{\E{\Wt^2}
=\frac{2(\lambda-\mu)}{\delta}V
-\frac{\lambda+\mu}{\delta}I
-K
+\frac{\lambda+\mu}{\delta}
+1,}
\label{dist:eq:ew2}
\end{equation}
with $\var{\Wt}=\E{\Wt^2}-V^2$. At $\mu=0$ and $t\to\infty$ it reduces to
$2\lambda^2/\delta^2+3\lambda/\delta+1$, the second moment of the geometric
burst-size law of \cref{dist:sec:burst-law} at success probability
$s=\delta/(\lambda+\delta)$.

Every formula above is a moment, and every one is quoted rather than derived
here. What produces them has not been written down yet.
